\chapter{Conclusion}

This dissertation has advanced the study of structural relaxation in glass-forming liquids by combining new experimental capability with a physically grounded framework for interpreting dynamics in thermodynamic space. The central result is a shift from describing relaxation along isolated thermodynamic paths to treating it as a continuous surface, $\tau(P,T)$, constrained by both experiment and physical boundary conditions.

On the experimental side, depolarized photon correlation spectroscopy (DPCS) was successfully extended into a diamond anvil cell environment, enabling direct measurement of structural $\alpha$-relaxation times at pressures approaching 2.5~GPa. This represents a significant expansion of the accessible thermodynamic space for optical probes of glassy dynamics, overcoming challenges associated with reduced scattering volume, heterodyne contamination, and photon-limited statistics. These measurements establish that DPCS remains a viable and powerful probe of structural relaxation under extreme conditions.

On the modeling side, relaxation dynamics were recast as a pressure--temperature surface, allowing fragility to be interpreted geometrically rather than as a purely isobaric quantity. By anchoring this surface to experimentally determined isochrones and incorporating physically motivated constraints, including a zero-mobility limit, the resulting framework avoids arbitrary extrapolation and provides a consistent description across broad thermodynamic ranges. Integration of high-pressure DPCS data with complementary dielectric and Brillouin measurements yields a unified representation of relaxation behavior.

Taken together, these results demonstrate that structural relaxation is inherently multi-dimensional, and that meaningful interpretation requires investigation across both temperature and pressure dimensions. The development of a thermodynamic surface model allows for the study of fragility as a function of both parameters and provides a pathway toward a more complete understanding of glass-forming systems.

Further refinement of the thermodynamic surface may be achieved through expansion of the TGP framework to multiple isochrones. By constraining the surface at several fixed relaxation times, rather than a single isochronal condition, the model can be more tightly anchored in the region surrounding the glass transition, where the most significant dynamical changes occur.

While Vogel--Tammann--Fulcher (VTF) behavior, and its modified forms along isothermal paths, appears to describe the dynamics of relatively simple liquids such as cumene with reasonable accuracy, it is well established that such models do not universally capture liquid dynamics across their full dynamic range. In contrast, our results indicate that the Andersson--Andersson formulation provides an exceptionally good description of the data along isochronal paths for both cumene and the hydrogen-bonded system glycerol---even through the transition associated with hydrogen-bond suppression. This observation raises the possibility that the Andersson--Andersson framework may serve not only as a high-pressure constraint, but as a more fundamental basis for constructing the full relaxation-time surface.

If such an approach proves broadly applicable, it invites deeper investigation into why this formulation succeeds under isochronal conditions. One possible interpretation is that isochronal paths, closely related to isochoric trajectories, more directly reflect the underlying physics governing structural relaxation. Future work may focus on testing this hypothesis and determining whether isochronal representations provide a more natural framework for describing glass-forming dynamics.